\begin{lemma}[Single-site geometry of the native torus]%
    \label{lem:peps_torus_singleton_geometry}
    \lean{TNLean.PEPS.torusContiguousRectangle_one_eq_singleton,
        TNLean.PEPS.torusGraph_compl_singleton_connected,
        TNLean.PEPS.torusGraph_singleton_connected,
        TNLean.PEPS.torusIncidentLeg_surjective,
        TNLean.PEPS.torusIncidentLeg_injective,
        TNLean.PEPS.torusIncidentLegEquiv,
        TNLean.PEPS.card_regionBoundaryEdge_torus_singleton,
        TNLean.PEPS.torusGraph_connected}
    \leanok
    \uses{def:peps_torus_geometry_translation,
        lem:peps_torus_rectangle_cut_connectivity}
    Suppose both periods of the native square torus are at least three.
    The graph is connected. For every vertex $v$, the subgraphs
    induced by $\{v\}$ and its complement are connected, and the
    boundary of $\{v\}$ has exactly four bonds. The native top, right,
    down, and left legs enumerate these incident bonds bijectively.
\end{lemma}

\begin{proof}\leanok
    A singleton is the coordinate rectangle of side lengths one.
    Its complement is connected by the unused row and column
    construction. The same rectangular connectivity argument applied
    to both full coordinate intervals gives connectivity of the torus.
    The four native legs exhaust the incident bonds and are distinct:
    horizontal and vertical bonds cannot agree, and opposite bonds
    have different origins when the corresponding period exceeds two.
    Every incident bond has just one endpoint at $v$, so these are
    precisely the crossing bonds of the singleton.
\end{proof}

\begin{theorem}[Single-site physical support of a regular torus PEPS]%
    \label{thm:peps_regular_torus_singleton_support}
    \lean{TNLean.PEPS.torusSingletonBoundaryEquiv,
        TNLean.PEPS.range_regionReducedDensity_torus_singleton_of_regular,
        TNLean.PEPS.rank_regionReducedDensity_torus_singleton_of_regular,
        TNLean.PEPS.stateCoeff_groupBondTensor_torus_ne_zero_of_regular,
        TNLean.PEPS.card_group_cube_le_physicalDimension_regular_torus,
        TNLean.PEPS.entropy_normalizedRegionReducedDensity_torus_singleton_of_regular}
    \leanok
    \uses{lem:peps_torus_singleton_geometry, def:peps_g_injective,
        def:peps_normalized_region_reduced_density}
    Let regular $G$-injective tensors of physical dimension $d$ be
    placed at the vertices of the native square torus, with both
    periods at least three. The tensors may vary with the vertex.
    Their closed vector is nonzero. For every vertex $v$,
    \[
        \operatorname{ran}\rho_{\{v\}}=\mathcal G_{\{v\}},
        \qquad
        \operatorname{rank}\rho_{\{v\}}=|G|^3\leq d.
    \]
    The normalized density satisfies
    \[
        0\leq S(\widehat\rho_{\{v\}})\leq3\log|G|,
        \qquad S_0(\widehat\rho_{\{v\}})=3\log|G|.
    \]
    No isometry or flat spectrum is assumed.
\end{theorem}

\begin{proof}\leanok
    \uses{thm:peps_regular_region_support,
        thm:peps_regular_region_dimension,
        thm:peps_regular_region_entropy_bound,
        lem:peps_torus_singleton_geometry}
    Apply the connected regular-cut statements to the singleton and
    its connected complement. Its four crossing bonds give three
    relative boundary labels. The density acts on a $d$-dimensional
    physical space, so its rank $|G|^3$ is at most $d$.
\end{proof}

\begin{theorem}[Local rigidity for equal regular torus states]%
    \label{thm:peps_regular_torus_local_rigidity}
    \lean{TNLean.PEPS.regionGroundSpace_torus_singleton_eq_of_sameState_regular,
        TNLean.PEPS.card_group_eq_of_sameState_regular_torus}
    \leanok
    \uses{lem:peps_torus_singleton_geometry, def:peps_g_injective,
        def:peps_region_ground_space}
    On a native square torus with both periods at least three, two
    regular $G$-injective descriptions of the same closed physical
    vector have the same single-site ground spaces at every vertex.
    If the two descriptions use regular representations of different
    finite groups $G$ and $H$, the equality of their closed vectors
    implies $|G|=|H|$. No assertion about the multiplication laws or
    individual bond gauges follows from this rank argument.
\end{theorem}

\begin{proof}\leanok
    \uses{thm:peps_same_state_regular_region_spaces,
        thm:peps_regular_group_order_rigidity,
        lem:peps_torus_singleton_geometry}
    Equal closed vectors give equal reduced densities. Their ranges
    recover the local ground spaces. The common single-site density
    has rank both $|G|^3$ and $|H|^3$, so the group orders agree.
\end{proof}

\begin{theorem}[Physical support of a bounded regular torus rectangle]%
    \label{thm:peps_regular_torus_rectangle_support}
    \lean{TNLean.PEPS.range_regionReducedDensity_torus_rectangle_of_regular}
    \leanok
    \uses{def:peps_g_injective, def:peps_region_ground_space,
        def:peps_region_reduced_density}
    Consider regular $G$-injective tensors on a native square torus
    with periods at least three. Let $R$ be a coordinate rectangle
    of positive side lengths, bounded by the coordinate intervals,
    with each side length strictly smaller than the corresponding
    torus period. Then $\operatorname{ran}\rho_R=\mathcal G_R$.
\end{theorem}

\begin{proof}\leanok
    \uses{thm:peps_regular_region_support}
    The rectangle is connected, and its complement is connected by
    paths along an unused row and an unused column. The torus is
    connected, so the exact support theorem applies.
\end{proof}
